\chapter{The accounting language}
\label{ch:language}

\section{Scope of the first language}

The historical AccDSL Core 0 proposal describes entities and books, accounts, observations, shared
unknowns, constraints, alternative interpretations, contract instances,
accounting rules and reporting mappings. It has pure expressions and a
restricted posting effect. Economic forecasting algorithms can live in typed
plugins; unrestricted numerical code does not need to be reimplemented in
the DSL. The historical grammar in Appendix~\ref{app:contracts} fixes the preserved
teaching examples. The proposed portable records and first company profile
below define the new authoring boundary; the first matched text frontend
supports the declared debt fragment.

A subsequent documentary subset authors the funded-year extension of that
company profile: identified claim allocations, explicit credit terms and a
dated event sequence. Text and typed records lower to the same body, which
drives the conditional reference calculation in
Section~\ref{sec:funded-year-authoring}. The imported annual projection still
supplies the other operating and accounting rules. This extends the matched
example without establishing a complete whole-company text grammar.

The next subset makes the operating inputs explicit as well.
Section~\ref{sec:company-input-authoring} starts from the historical opening,
then derives purchases, sales, claims, expenses and closing statements from
authored quantities, prices, asset cohorts and contract terms. It imports no
forecast journals. This makes a price or service-date change traceable through
the accounts, while retaining the example's stated method restrictions.

Identified operations can then repeat within a year and carry their unpaid
claims and asset cohorts into the next year. The conditional path in
Section~\ref{sec:company-operation-paths} follows Northstar through three
annual periods, including revised assumptions and an interrupted year.
Its typed records specify this composition. The bounded text frontend in
Section~\ref{sec:company-path-authoring} preserves those records exactly,
including named settlement allocations and incomplete paths.
Section~\ref{sec:company-family-language} adds named unweighted scenarios and
pure typed expressions for the declared collection policy.
Section~\ref{sec:company-module-language} preserves their portable identities
and specifies the additional state needed by a shared-book return method.
Section~\ref{sec:company-return-reference} gives its matched authoring
and finite shared-book execution reference. Section~\ref{sec:company-source-state}
connects explicit company method selection to source-constrained partial state.
Section~\ref{sec:statement-context-relations} extends the retained source
representation to all statement observations and their distinct relation roles.
Section~\ref{sec:joint-source-measurements} then tests selected balance relations
jointly under explicit rounding and identity hypotheses.
Section~\ref{sec:temporal-statement-network} preserves opening, movement and
closing times while extending the selected system across the statements.
Section~\ref{sec:conditional-history-families} gives a finite conditional
history family with symbolic openings and explicit measurement attachment.
Section~\ref{sec:typed-economic-history} adds typed economic quantities,
constraints at every event and matched declarations for conditional families.
Section~\ref{sec:typed-history-observation-attachment} connects cash and
quantity observations to the same complete economic history.
Section~\ref{sec:conditional-method-families} connects unknown monetary
amounts to selected method prerequisites without choosing a single history.
Section~\ref{sec:native-conditional-sequence} carries every retained return
branch through settlement and sale, preserving funding and intermediate observations.
Section~\ref{sec:family-outcome-queries} asks which financial outcomes are possible
or necessary, and how exact finite bounds retain those alternatives.
Section~\ref{sec:simulator-evidence-archive} preserves complete definitions and
results across restart while separating stored evidence from verified execution.
Section~\ref{sec:lease-shared-book} composes heterogeneous lease rights in
one company book with explicit account allocation and funding conditions.
Section~\ref{sec:lease-allocation-family} preserves the allocations left
unidentified by aggregate balances and their different complete or interrupted outcomes.
Section~\ref{sec:henkel-lease-measurements} tests what actual rounded lease
disclosures can constrain before individual contract state is known.
Section~\ref{sec:shared-return-lease-book} then distinguishes a common caller
from compatible state and follows return and lease events through one declared cash balance.
Section~\ref{sec:recorded-shared-statements} reconciles the resulting recorded statements
while preserving unresolved tax assessment and incomplete fiscal coverage.
Section~\ref{sec:shared-book-authoring} authors the complete opening and mixed
event sequence through matched declarations and typed records.
Section~\ref{sec:shared-tax-information} separates supplied tax declarations
from executable tax rules and explains the remaining mixed-book incompatibility.
Section~\ref{sec:method-requirement-language} makes selected method prerequisites
reusable across these profiles, with matched typed and text declarations.
Section~\ref{sec:requirement-successors} connects exact input binding to
dimensional products and guarded applicability, with explicit version selection.
Section~\ref{sec:collection-state-requirements} adds explicit collection coverage,
identity order, signed totals and state-domain findings before monetary effects.
Section~\ref{sec:typed-return-transitions} carries the same example through
typed calculations, atomic state updates and exact source-bound replay.
General functions,
endogenous decisions, source admission and broader accounting profiles remain
language tasks.

The historical Core 0 compiler has six ordered stages: parsing with source locations; name and
module resolution; type and temporal-context checking; lowering to the IR;
constraint and rule-applicability analysis; and execution-plan construction.
The compiler never silently supplies an economic assumption to make a program
compile. A program can be valid as a historical representation while requiring
additional assumptions before a particular simulation can run.
The proposed portable compilation sequence is specified separately in
Section~\ref{sec:portable-authoring}.

The language's meaning is defined by its IR and evaluator, rather than by the
surface spelling alone. Formatting changes cannot change a model's semantic
hash. A changed source span or reviewer decision changes evidence provenance
even if the numerical model remains the same. The system therefore maintains
separate evidence, semantic-model and execution identifiers.

For the preserved Core 0 examples, the complete grammar is supplied as
\code{spec/core-0.lark}; the original design directory retains
\code{spec/core-ir.schema.json} and \code{spec/core-contract.md}. They define reference resolution, expression arity,
type/context checks, operations, failures and identity. Structural schema
validation alone does not establish those semantic properties. ASCII
identifiers and UTF-8 strings are used in this revision; currency suffixes,
dates and decimal literals have explicit lexical rules. The grammar is checked
by parsing the complete demonstration file, not by claiming that its parser
also implements a compiler or solver.

